\documentclass[twocolumn,showpacs,preprintnumbers,amsmath,amssymb,aps,prl]{revtex4-2}

\usepackage{graphicx}
\usepackage{dcolumn}
\usepackage{bm}
\usepackage{hyperref}
\usepackage{booktabs}

\begin{document}

\title{Supplemental Material:\\
Spatial Scrambling Profiles and System-Size Scaling of $\Omega$\\
across Dynamical Regimes}

\author{Erick Francisco Perez Eugenio}
\email{erick.fpe79@gmail.com}
\affiliation{Independent Researcher, Le\'on, Guanajuato, Mexico,
ORCID: 0009-0006-3228-4847}

\date{March 2026}

\begin{abstract}
This Supplemental Material contains: (S1) complete $\Omega_{\rm late}(r)$
profiles for all five models and all system sizes; (S2) full cross-platform
hardware data tables; (S3) fit diagnostics and AIC computation; (S4) circuit and
simulation parameters.

All numerical values in this document are extracted programmatically from the
authoritative data files (\texttt{paper1b\_master\_5models.json},
\texttt{aqt\_results.json}, \texttt{rigetti\_results.json}) and verified against
independent recomputation. The extraction script (\texttt{generate\_figures.py})
is included in the repository.
\end{abstract}

\maketitle

% ============================================================
\section{Complete Spatial Profiles $\Omega_{\rm late}(r)$}
% ============================================================

Tables~\ref{tab:kichaotic1}--\ref{tab:syk} give the complete
$\Omega_{\rm late}(r)$ profiles for each model. All values from exact statevector
simulation with $D_{\max} = 5N$, dense trace method ($N \leq 12$) or typicality
with 10 random states ($N \geq 14$, KI chaotic only).

\begin{table}[h]
\caption{KI chaotic $\Omega_{\rm late}(r)$ profiles, $r = 1$--$7$. Dash (\mbox{---})
indicates $r \geq N$.}
\label{tab:kichaotic1}
\begin{ruledtabular}
\begin{tabular}{cccccccc}
$N$ & $r{=}1$ & $r{=}2$ & $r{=}3$ & $r{=}4$ & $r{=}5$ & $r{=}6$ & $r{=}7$ \\
\hline
4  & 0.280 & 0.326 & 0.331 & ---   & ---   & ---   & ---   \\
6  & 0.215 & 0.252 & 0.274 & 0.275 & 0.402 & ---   & ---   \\
8  & 0.159 & 0.191 & 0.210 & 0.222 & 0.244 & 0.262 & 0.393 \\
10 & 0.135 & 0.162 & 0.180 & 0.186 & 0.197 & 0.210 & 0.227 \\
12 & 0.119 & 0.140 & 0.156 & 0.159 & 0.168 & 0.176 & 0.185 \\
\end{tabular}
\end{ruledtabular}
\end{table}

\begin{table}[h]
\caption{KI chaotic $\Omega_{\rm late}(r)$ profiles, $r = 8$--$11$. Only $N=10$
and $N=12$ have data in this range.}
\label{tab:kichaotic2}
\begin{ruledtabular}
\begin{tabular}{ccccc}
$N$ & $r{=}8$ & $r{=}9$ & $r{=}10$ & $r{=}11$ \\
\hline
10 & 0.247 & 0.376 & ---   & ---   \\
12 & 0.196 & 0.212 & 0.234 & 0.360 \\
\end{tabular}
\end{ruledtabular}
Note: the last distance ($r = N-1$) in OBC consistently shows a jump in
$\Omega_{\rm late}$ (e.g., $N=6$: $r=4$ gives 0.275 but $r=5$ gives 0.402;
$N=12$: $r=10$ gives 0.234 but $r=11$ gives 0.360). This boundary effect, where
the far-end operator is partially shielded from scrambling by its proximity to
the open boundary, is present at all system sizes and should not be confused with
the monotonic increase due to butterfly velocity at intermediate distances.
\end{table}

\begin{table}[h]
\caption{KI mixing $\Omega_{\rm late}(r)$ profiles, $r = 1$--$7$. The U-shaped
pattern (minimum at interior, maxima at boundaries) persists across all $N$ but
becomes asymmetric for $N \geq 10$: the far-boundary values are lower than the
near-boundary values.}
\label{tab:kimixing1}
\begin{ruledtabular}
\begin{tabular}{cccccccc}
$N$ & $r{=}1$ & $r{=}2$ & $r{=}3$ & $r{=}4$ & $r{=}5$ & $r{=}6$ & $r{=}7$ \\
\hline
4  & 0.533 & 0.461 & 0.820 & ---   & ---   & ---   & ---   \\
6  & 0.865 & 0.599 & 0.472 & 0.632 & 0.897 & ---   & ---   \\
8  & 0.953 & 0.824 & 0.618 & 0.536 & 0.579 & 0.738 & 0.916 \\
10 & 0.977 & 0.922 & 0.786 & 0.622 & 0.573 & 0.586 & 0.660 \\
12 & 0.985 & 0.956 & 0.895 & 0.764 & 0.630 & 0.588 & 0.593 \\
\end{tabular}
\end{ruledtabular}
\end{table}

\begin{table}[h]
\caption{KI mixing $\Omega_{\rm late}(r)$ profiles, $r = 8$--$11$. Asymmetry
visible: compare $r=9$ at $N=10$ (0.932) with $r=1$ (0.977); compare $r=11$ at
$N=12$ (0.952) with $r=1$ (0.985).}
\label{tab:kimixing2}
\begin{ruledtabular}
\begin{tabular}{ccccc}
$N$ & $r{=}8$ & $r{=}9$ & $r{=}10$ & $r{=}11$ \\
\hline
10 & 0.811 & 0.932 & ---   & ---   \\
12 & 0.635 & 0.725 & 0.853 & 0.952 \\
\end{tabular}
\end{ruledtabular}
\end{table}

\begin{table}[h]
\caption{SYK $\Omega_{\rm late}(r)$ profiles. Disorder average over 10 seeds.
Flat profile to within $\lesssim 10^{-3}$ at each $N$, confirming absence of
spatial structure from all-to-all connectivity.}
\label{tab:syk}
\begin{ruledtabular}
\begin{tabular}{cccccccccc}
$N$ & $r{=}1$ & $r{=}2$ & $r{=}3$ & $r{=}4$ & $r{=}5$ & $r{=}6$ & $r{=}7$
    & $r{=}8$ & $r{=}9$ \\
\hline
4  & 0.072 & 0.072 & 0.071 & ---   & ---   & ---   & ---   & ---   & ---   \\
6  & 0.016 & 0.018 & 0.019 & 0.020 & 0.019 & ---   & ---   & ---   & ---   \\
8  & 0.005 & 0.005 & 0.004 & 0.005 & 0.005 & 0.005 & 0.005 & ---   & ---   \\
10 & 0.001 & 0.001 & 0.001 & 0.001 & 0.001 & 0.001 & 0.001 & 0.001 & 0.001 \\
\end{tabular}
\end{ruledtabular}
\end{table}

Integrable and dual-unitary models give $\Omega_{\rm late}(r) = 1.000$ to machine
precision ($< 10^{-13}$) for all $N$ and $r$; these values are omitted for
brevity.

% ============================================================
\section{Cross-Platform Hardware Data}
% ============================================================

Tables~\ref{tab:aqt} and~\ref{tab:rigetti} give the full $\langle Z_0 \rangle$
measurements for $N=4$ across all circuit types and depths. Each row corresponds
to one circuit: the magnetization (mag) circuit measures $\langle Z_0 \rangle$
after forward evolution $U_F^d |\psi_0\rangle$; the echo circuit measures
$\langle Z_0 \rangle$ after the full echo $U_F^{-d} Z_r U_F^d |\psi_0\rangle$
with $r=1$.

\begin{table}[h]
\caption{AQT IBEX Q1: $\langle Z_0 \rangle$ measurements vs.\ exact simulation.
$N=4$, 100 shots per circuit ($\sigma \approx 0.10$). Values are qubit-0
magnetization.}
\label{tab:aqt}
\begin{ruledtabular}
\begin{tabular}{llccrr}
Model & Type & $d$ & Exact & AQT & $|\Delta|$ \\
\hline
Integrable & mag  & 1 & $+0.540$ & $+0.500$ & 0.040 \\
Integrable & echo & 1 & $-0.416$ & $-0.280$ & 0.136 \\
Integrable & mag  & 3 & $-0.990$ & $-1.000$ & 0.010 \\
Integrable & echo & 3 & $+0.960$ & $+0.880$ & 0.080 \\
Integrable & mag  & 5 & $+0.284$ & $+0.340$ & 0.056 \\
Integrable & echo & 5 & $-0.839$ & $-0.820$ & 0.019 \\
KI chaotic & mag  & 1 & $+0.540$ & $+0.420$ & 0.120 \\
KI chaotic & echo & 1 & $-0.416$ & $-0.440$ & 0.024 \\
KI chaotic & mag  & 3 & $+0.575$ & $+0.560$ & \textbf{0.015} \\
KI chaotic & echo & 3 & $+0.126$ & $-0.140$ & 0.266 \\
KI chaotic & mag  & 5 & $+0.550$ & $+0.420$ & 0.130 \\
KI chaotic & echo & 5 & $+0.074$ & $+0.000$ & 0.074 \\
\end{tabular}
\end{ruledtabular}
\end{table}

\begin{table}[h]
\caption{Rigetti Ankaa-3: $\langle Z_0 \rangle$ measurements vs.\ exact
simulation. $N=4$, 2048 shots ($\sigma < 0.02$). Values are qubit-0
magnetization.}
\label{tab:rigetti}
\begin{ruledtabular}
\begin{tabular}{llccrr}
Model & Type & $d$ & Exact & Rigetti & $|\Delta|$ \\
\hline
Integrable & mag  & 1 & $+0.540$ & $+0.517$ & 0.024 \\
Integrable & echo & 1 & $-0.416$ & $-0.376$ & 0.040 \\
Integrable & mag  & 3 & $-0.990$ & $-0.867$ & 0.123 \\
Integrable & echo & 3 & $+0.960$ & $+0.924$ & 0.036 \\
Integrable & mag  & 5 & $+0.284$ & $+0.248$ & 0.036 \\
Integrable & echo & 5 & $-0.839$ & $-0.731$ & 0.108 \\
KI chaotic & mag  & 1 & $+0.540$ & $+0.546$ & 0.006 \\
KI chaotic & echo & 1 & $-0.416$ & $-0.077$ & 0.339 \\
KI chaotic & mag  & 3 & $+0.575$ & $+0.356$ & \textbf{0.219} \\
KI chaotic & echo & 3 & $+0.126$ & $-0.125$ & 0.251 \\
KI chaotic & mag  & 5 & $+0.550$ & $+0.112$ & 0.437 \\
KI chaotic & echo & 5 & $+0.074$ & $-0.032$ & 0.106 \\
\end{tabular}
\end{ruledtabular}
\end{table}

The key comparison is the bolded entries: KI chaotic $d=3$ magnetization has
$|\Delta| = 0.015$ (AQT) vs.\ $|\Delta| = 0.219$ (Rigetti), a $15\times$
improvement. Echo circuits are more demanding (double the effective circuit depth)
and both platforms show larger errors for echo at $d \geq 3$.

\textit{Note on IBM data.}
IBM ibm\_marrakesh data (429 circuits, Conjunto~A parameters) are reported in
full in Paper~I~\cite{PerezEugenio2026a}. They are included here only in
aggregated form (Table~IV of the main text) for protocol comparison. IBM data
use Conjunto~A; Rigetti and AQT use Conjunto~B. The cross-platform comparison
validates the measurement protocol, not point-by-point physical equivalence.

\textit{Note on echo q1--q3 in AQT integrable.}
For all three depths, AQT IBEX Q1 measures $\langle Z \rangle = 1.000$ exactly on
qubits 1, 2, and~3 in the integrable echo circuit. With $J=0$ (no entangling
gates), the perturbation $Z_1$ at site~1 does not propagate to distant qubits;
the echo should return them to the initial state $|0\rangle$ with
$\langle Z \rangle = 1$ exactly. That 100 out of 100 shots give $|0\rangle$ on
all three qubits at all depths confirms exceptional single-qubit fidelity of the
ion trap.

% ============================================================
\section{Fit Diagnostics}
% ============================================================

\subsection{AIC Computation}

The Akaike Information Criterion is ${\rm AIC} = 2k - 2\ln\hat{L}$, where $k$ is
the number of free parameters and $\hat{L}$ is the maximized likelihood. For
least-squares fits to $n$ data points with residual sum of squares RSS:
\begin{equation}
  {\rm AIC} = n\ln({\rm RSS}/n) + 2k.
\end{equation}
Both algebraic (Hip.~A) and power-law (Hip.~C) have $k=2$ free parameters, so
$\Delta{\rm AIC}$ depends only on the ratio of RSS values. With
${\rm RSS}_A = 3.21\times10^{-5}$ and ${\rm RSS}_C = 4.67\times10^{-5}$ for
$n=8$:
\begin{equation}
  \Delta{\rm AIC} = 8\ln\!\left(\frac{3.21}{4.67}\right) = -3.00.
\end{equation}
By convention, $|\Delta{\rm AIC}| > 2$ constitutes positive evidence and
$|\Delta{\rm AIC}| > 6$ constitutes strong evidence. The exact RSS values from
the numerical fit (before rounding) give $\Delta{\rm AIC} = -3.30$, reported
throughout the main text; the rounded values quoted above yield $-3.00$ as an
approximate illustration of the computation. The value $-3.30$ represents
moderate positive evidence in favor of Hip.~A.

\subsection{Exponential Exclusion (Hip.~B)}

The exponential fit $\Omega = e^{-\gamma N}$ was tested on the same 8-point
series. Best-fit $\gamma = 0.053$, ${\rm RSS}_B = 2.48\times10^{-3}$, giving
$R^2 = 0.967$ and $\Delta{\rm AIC} = +19.1$ relative to power-law. This provides
very strong empirical evidence against the exponential form for KI chaotic. For
systems with energy conservation, the Huang-Brand\~ao-Zhang
theorem~\cite{HuangFiniteSize2019} provides additional theoretical support;
KI chaotic is Floquet (no energy conservation), so its exclusion rests on
empirical grounds alone.

\subsection{$\beta$-Function Finite-Difference Table}

\begin{table}[h]
\caption{Effective $\beta$-function from 8-point KI chaotic series, central
differences. $\Omega$ column shows the value at the central point;
$\beta_{\rm eff}$ is $(\Omega_{i+1}-\Omega_{i-1})/(\ln N_{i+1} - \ln N_{i-1})$.}
\label{tab:beta}
\begin{ruledtabular}
\begin{tabular}{ccc}
$\Omega$ & $\beta_{\rm eff}$ & $-c\,\Omega(1-\Omega)$ \\
\hline
0.215 & $-0.175$ & $-0.167$ \\
0.159 & $-0.157$ & $-0.132$ \\
0.135 & $-0.098$ & $-0.116$ \\
0.119 & $-0.096$ & $-0.104$ \\
0.103 & $-0.098$ & $-0.091$ \\
0.091 & $-0.085$ & $-0.082$ \\
\end{tabular}
\end{ruledtabular}
\end{table}

The data show that $\beta_{\rm eff}$ is systematically more negative than
$-c\,\Omega(1-\Omega)$ at large $\Omega$ (small $N$), which is expected from
finite-size effects in the pre-asymptotic regime. At small $\Omega$ (large $N$),
the two converge, consistent with the algebraic fit becoming more accurate as
finite-size corrections diminish.

\begin{figure}[h]
  \includegraphics[width=\columnwidth]{fig3_beta_function}
  \caption{Effective $\beta$-function
  $\beta_{\rm eff} = \Delta\Omega/\Delta(\ln N)$ vs.\ $\Omega_{\rm late}$ for
  KI chaotic (circles). Algebraic fit (solid, $\beta = -c\,\Omega(1-\Omega)$,
  $c = 0.990$) and power-law fit (dashed, $\beta = -\alpha\,\Omega$,
  $\alpha = 0.810$) are statistically indistinguishable in the data range
  $\Omega \in [0.08, 0.28]$.}
  \label{fig:beta}
\end{figure}

% ============================================================
\section{Circuit and Simulation Parameters}
% ============================================================

\subsection{Kicked Ising Circuit (Conjunto~B)}

The Floquet operator is:
\begin{equation}
  U_F = e^{-ih\sum_{i=0}^{N-1} X_i}\,
        e^{-ib\sum_{i=0}^{N-1} Y_i}\,
        e^{-iJ\sum_{i=0}^{N-2} Z_i Z_{i+1}},
\end{equation}
with $J=1.0$, $h=0.5$, $b=0.5$ (KI chaotic), $J=0.1$ (KI mixing),
$J=\pi/4$, $h=\pi/4$, $b=0$ (dual-unitary), $J=0$, $b=0$ (integrable). Open
boundary conditions; no $Z_i Z_{i+1}$ term for $i=N-1$.

\subsection{Echo Protocol}

Initial state: $|\psi_0\rangle = |0\rangle^{\otimes N}$ (all qubits in $|0\rangle$
state). Circuit structure per depth $d$, scrambled operator at site $r$:
\begin{equation}
  |\phi_r(d)\rangle = U_F^{-d}\,Z_r\,U_F^d\,|\psi_0\rangle.
\end{equation}
Observable: $\langle Z_0 \rangle = \langle \phi_r(d) | Z_0 | \phi_r(d) \rangle$.
Normalization: $C(0) = \langle\psi_0|Z_0|\psi_0\rangle = 1$ (all qubits in
$|0\rangle$).

\subsection{$D_{\max}$ and $\Omega_{\rm late}$ Computation}

For each $(N, \text{model})$ pair, $D_{\max} = 5N$. The late-time average uses
depths $d \geq D_{\max}/2$, i.e., the second half of the circuit depth range.
For $N=4$: depths $d = 10,\ldots,20$ contribute. The dense trace method computes
the exact statevector at each depth; the typicality method approximates the trace
by averaging over 10 random initial states sampled from the Haar measure on the
$2^N$-dimensional Hilbert space.

\subsection{SYK Model}

The SYK-inspired Floquet operator uses random all-to-all ZZ couplings:
$H_{\rm SYK} = \sum_{i<j} J_{ij} Z_i Z_j$ with $J_{ij} \sim \mathcal{U}(0.5,
1.5)$ independently for each realization. The Floquet operator is
$U_{\rm SYK} = e^{-iH_{\rm SYK}}$. Disorder averaging uses 10 independent
coupling realizations for all $N \leq 12$.

\subsection{Hardware Circuit Transpilation}

All circuits were transpiled using Qiskit 2.3.0 (IBM), PyQuil (Rigetti via AWS
Braket), and the AQT SDK (AQT via AWS Braket), with optimization level~3. The
echo protocol was verified with integrable circuits on all platforms before
running chaotic models.

% ============================================================
\section{Repository Contents}
% ============================================================

The data repository~\cite{PerezEugenio2026a} contains:

\textit{Data files:}
\begin{itemize}
  \item \texttt{paper1b\_master\_5models.json} --- authoritative dataset, 5 models,
    $N=4$--$12$, $D_{\max}=5N$, dense trace
  \item \texttt{ki\_chaotic\_N14\_18.json} --- KI chaotic typicality results,
    $N=14,16,18$
  \item \texttt{syk\_n12\_results.json} --- SYK $N=12$, 10 seeds, $D_{\max}=60$
  \item \texttt{aqt\_results.json} --- 12 AQT circuits, raw magnetizations (all
    4 qubits)
  \item \texttt{aqt\_analysis\_final.json} --- 48 AQT measurements with
    exact-simulation diffs
  \item \texttt{rigetti\_results.json} --- 26 Rigetti circuits, raw
  \item \texttt{rigetti\_analysis\_final.json} --- Rigetti analysis
\end{itemize}

\textit{Scripts:}
\begin{itemize}
  \item \texttt{omega\_simulation.py} --- exact simulation (dense trace and
    typicality)
  \item \texttt{scaling\_fits.py} --- algebraic/power-law fits, $\Delta$AIC,
    $\beta_{\rm eff}$ computation
  \item \texttt{cross\_platform\_analysis.py} --- hardware comparison
  \item \texttt{generate\_figures.py} --- reproduces all figures from JSON files
\end{itemize}

\textit{Documents:}
\begin{itemize}
  \item \texttt{paper1b\_main.tex} --- main manuscript
  \item \texttt{paper1b\_supplemental.tex} --- this document
  \item \texttt{kaelion\_v8.bib} --- bibliography (v8, 58 entries)
\end{itemize}

\bibliography{kaelion_v8}

\end{document}
